%% ============================================================
%% nro/equipe.tex — o que todo relatório sabe da equipe
%%
%% Mantido pela DIRETORIA. O membro não mexe aqui: se o professor que o
%% avalia não é o da equipe, ou se a norma do colegiado dele é outra,
%% ele repete \equipe{…} no preâmbulo do relatório só com o que muda.
%%
%% Antes de publicar o template, a Diretoria preenche o que está vazio
%% abaixo. Campo vazio vira pendência em todos os relatórios — de
%% propósito: sem o professor responsável e sem o registro no SIEX, o
%% relatório não se sustenta diante do Colegiado.
%% ============================================================

%% A série no controle de arquivos (NRO-PUB-001). É um template de
%% registros: cada relatório é um PN, que o membro cria em Arquivos,
%% no SOMA, e sobe lá como revisão para o grupo revisor aprovar.
%% O código é uma PROPOSTA: confira em Arquivos que o SN 020 do Pessoal
%% está livre, crie a série com a estrutura "Template → registros" e,
%% se o número mudar, troque-o aqui.
\serie{
  codigo  = NRO-PES-020,
  titulo  = Relatório de atividades para aproveitamento de AACC,
  emissor = Departamento de Pessoal,
  revisao = A,
}

\equipe{
  nome         = NeuroDynamics PD\&I,
  sigla        = NRO,
  %% como a frase o escreve: "realizadas no …"
  laboratorios = {Laboratório de Engenharia Biomédica e no Laboratório de Bioengenharia (LABBIO) da Escola de Engenharia da UFMG},
  site         = neurodynamics.dev,
  validacao    = auth.neurodynamics.dev,

  %% o professor responsável pela equipe: quem avalia e dá a nota.
  %% Com o título, como deve aparecer na assinatura (Prof. Dr. …,
  %% Profa. Dra. …); o artigo (o ou a) faz o texto dizer "do" ou "da".
  professor             = {},
  artigoprofessor       = o,
  departamentoprofessor = {},   % ex.: Departamento de Engenharia Mecânica
  emailprofessor        = {},

  %% quem dirige a equipe e assina o parecer da equipe
  dirigente      = {},
  cargodirigente = {},          % ex.: Diretor-geral

  %% a ação de extensão, como está no SIEX/UFMG
  siextitulo        = {},
  siexnumero        = {},
  siexcoordenador   = {},        % com o título: Prof. Dr. …
  artigocoordenador = o,
  siexvigencia      = {},        % ex.: 2025 a 2027

  %% a norma que rege o aproveitamento; o membro de outra unidade troca
  %% no próprio relatório
  norma = {Resolução da Congregação da Escola de Engenharia da UFMG n\textordmasculine~02/2019, de 31 de maio de 2019},

  %% as contas: 1 crédito = 15 horas na UFMG; e a média semanal a partir
  %% da qual o relatório precisa justificar as horas
  horasporcredito = 15,
  mediamaxima     = 20,
}

\documento{
  cidade = Belo Horizonte,
}
